\section{Design Philosophy}

The pipeline is structured as a graph-level proofreader: the input segmentation
volume is reduced to a set of fragment skeletons in the first stage, and every
subsequent decision (graph construction, candidate scoring, rule-based
validation, assembly) operates on skeleton geometry and fragment-level summary
statistics --- never on raw voxels. The voxel segmentation is the early input
and the late rendering target; the computational substrate for split detection
is the skeleton. This places the pipeline in the same tradition as NEURD
\cite{celii2025neurd}, Joyce et al.~\cite{joyce2023mesh}, and
ConnectomeBench~\cite{brown2025connectomebench}, which likewise reason over
mesh or skeleton representations rather than image patches.

Within this overall design, three principles govern the decision logic. First,
\textbf{conservative correctness}: a merge that should not happen is worse
than a merge proposal that is correctly rejected, so the pipeline never forces
an ambiguous decision; uncertainty is output explicitly as a third outcome
(AMBIGUOUS). Second, \textbf{modularity}: each stage can be run independently,
with intermediate results saved to disk and re-used across parameter sweeps.
Third, \textbf{reproducibility}: all parameters are controlled through a
single YAML configuration file, and random seeds are fixed, so the same
config always produces the same output.

The six-stage pipeline is shown in Figure~\ref{fig:pipeline}. The remainder
of this chapter describes each stage.

\begin{figure}[H]
\centering
\begin{tikzpicture}[
  node distance = 4.5mm,
  every node/.style  = {align=center},
  %--- node styles ---
  io/.style = {
    draw, ellipse,
    minimum width = 62mm, minimum height = 9mm,
    fill = gray!12, draw = gray!55, line width = 0.65pt,
    font = \small\bfseries,
  },
  stage/.style = {
    draw, rounded corners = 5pt,
    minimum width = 90mm,
    inner xsep = 10pt, inner ysep = 7pt,
    fill = blue!7, draw = blue!35!black, line width = 0.65pt,
  },
  key/.style = {stage, fill = orange!7, draw = orange!40!black},
  out/.style = {stage, fill = green!7, draw = green!45!black},
  %--- arrow style ---
  arr/.style = {
    -{Stealth[length=5pt, width=4pt]},
    line width = 0.9pt, color = gray!60!black,
  },
]

%% Input node
\node[io] (inp) {Segmentation Volume};

%% Stage 1
\node[stage, below=of inp] (s1) {
  {\small\bfseries [1]\enspace Fragment Extraction}\\[2pt]
  {\footnotesize\color{gray!55!black}
    Connected-component extraction $\cdot$ TEASAR skeletonisation $\cdot$ chunk stitching}
};

%% Stage 2
\node[stage, below=of s1] (s2) {
  {\small\bfseries [2]\enspace Graph Construction}\\[2pt]
  {\footnotesize\color{gray!55!black}
    Skeleton-node KD-tree $\cdot$ optional long-range endpoint pass}
};

%% Stage 3
\node[stage, below=of s2] (s3) {
  {\small\bfseries [3]\enspace Candidate Generation}\\[2pt]
  {\footnotesize\color{gray!55!black}
    Composite scoring $\cdot$ distance-conditioned weights}
};

%% Stage 4 — highlighted as the key decision stage
\node[key, below=of s3] (s4) {
  {\small\bfseries [4]\enspace Conservative Validation}\\[2pt]
  {\footnotesize\color{gray!55!black}
    7 configurable rules $\cdot$ three outcomes: \textbf{accept} / \textbf{reject} / \textbf{ambiguous}}
};

%% Stage 5
\node[stage, below=of s4] (s5) {
  {\small\bfseries [5]\enspace Assembly}\\[2pt]
  {\footnotesize\color{gray!55!black}
    Connected-component merging $\cdot$ cycle detection $\cdot$ topology checks}
};

%% Stage 6 — output
\node[out, below=of s5] (s6) {
  {\small\bfseries [6]\enspace Export}\\[2pt]
  {\footnotesize\color{gray!55!black}
    CSV $\cdot$ GraphML $\cdot$ SWC $\cdot$ Neuroglancer annotations}
};

%% Arrows
\draw[arr] (inp) -- (s1);
\draw[arr] (s1)  -- (s2);
\draw[arr] (s2)  -- (s3);
\draw[arr] (s3)  -- (s4);
\draw[arr] (s4)  -- (s5);
\draw[arr] (s5)  -- (s6);

\end{tikzpicture}
\caption{Overview of the six-stage post-segmentation pipeline. Stage~4
(Conservative Validation, shaded orange) is the primary decision point; it
returns one of three explicit outcomes so that uncertainty is never silently
discarded.}
\label{fig:pipeline}
\end{figure}

\section{Stage 1: Fragment Extraction}

\paragraph{Connected component extraction.}
The input is a voxel-wise label volume (HDF5 format). Because full volumes
may exceed available memory ($699^3$ voxels at 64 bits $\approx$ 2.7~GB), the
pipeline reads the volume in configurable chunks with overlap, extracts
connected components per chunk, and stitches cross-boundary fragments using
label agreement in overlap zones.

Each resulting \texttt{Fragment} object stores its integer label identifier,
voxel count, axis-aligned bounding box, centroid (in physical nm coordinates),
a list of endpoint coordinates, an optional TEASAR skeleton, and a boundary flag
indicating whether the fragment touches a chunk boundary.

\paragraph{Skeletonization.}
Fragments above a minimum voxel threshold (100 voxels by default, configurable)
are skeletonized using the \texttt{kimimaro} TEASAR implementation. The TEASAR
invalidation radius parameter \texttt{invalidation\_d0} controls skeleton
branch density; we use \texttt{invalidation\_d0 = 10} for XPRESS.

For very large fragments (above a configurable \texttt{max\_skeleton\_voxels}
threshold, set to 500{,}000 voxels for XPRESS), TEASAR runtime scales
non-linearly and can require $>10$ minutes per fragment. For these, the pipeline
falls back to a PCA-based bounding-box endpoint: the fragment's bounding box
is projected onto its first principal component, and the two extreme corners
of the projected box are used as pseudo-endpoints.

\paragraph{Fragment stitching.}
The stitching module uses a union-find data structure to merge fragment IDs
across chunk boundaries, relabeling fragments globally after processing all
chunks. On the XPRESS training volume, 14{,}591 per-chunk fragments stitch
into 11{,}208 global fragments; on the held-out validation volume the same
procedure yields 12{,}935 global fragments.

\section{Stage 2: Graph Construction}

Graph construction determines which fragment pairs are considered as merge
candidates. This stage had the single largest impact on recall in our
experiments.

\paragraph{Endpoint graph (baseline).}
The naive approach indexes only TEASAR degree-1 endpoints (skeleton leaf nodes)
in a KD-tree and connects pairs within a distance threshold
(\texttt{max\_distance\_nm}). This correctly handles splits at the tips of axons
but cannot represent splits in the interior of long axons, where the split point
is surrounded by skeleton interior nodes, not endpoints.

\paragraph{Skeleton-node graph (proposed method).}
The \texttt{skeleton\_node} construction method indexes \emph{every} skeleton
node (leaf, branch, and interior) from every fragment in a single global
KD-tree. All node-to-node pairs within \texttt{max\_distance\_nm} receive an
edge. This exposes interior splits at any point along an axon.

On XPRESS, switching from the endpoint graph (at 300-voxel crop resolution) to
the skeleton-node graph (at full-volume 699-voxel resolution) raised
ground-truth coverage from 23\% to 79.8\%, a 56.8-percentage-point gain in one
change.

\paragraph{Long-range endpoint pass.}
A supplemental second pass queries only degree-1 endpoints at a larger radius
(\texttt{max\_endpoint\_search\_nm} = 2000~nm) using the same KD-tree.
This adds edges for genuine segmentation gaps that are wider than the skeleton-node
radius (500~nm), targeting true splits where the two fragments' interior skeleton
nodes are not close but their free endpoints point toward each other across the gap.

\section{Stage 3: Candidate Generation and Scoring}

\paragraph{Composite score.}
For each graph edge $(A, B)$, the candidate generator identifies the closest
endpoint pair using the skeleton node positions, then computes four normalized
sub-scores:

\begin{description}
  \item[Proximity score] $s_\text{prox} = \exp(-d / d_\text{ref})$, where $d$
        is the gap distance and $d_\text{ref}$ is a configurable reference
        distance (600~nm for XPRESS). Decays to zero as distance grows.

  \item[Alignment score] $s_\text{align} = (\hat{t}_A \cdot \hat{t}_B + 1) / 2$,
        where $\hat{t}_A$ and $\hat{t}_B$ are the unit tangent vectors at the
        endpoints of fragments $A$ and $B$, estimated from the local skeleton
        geometry. A score of 1.0 means the fragments point directly toward each
        other; 0.0 means they are antiparallel.

  \item[Continuity score] $s_\text{cont} = \max(0,\; 1 - \kappa / \kappa_\text{max})$,
        where $\kappa$ is the angular bend at the proposed junction and
        $\kappa_\text{max} = 180°$. Penalizes sharp turns.

  \item[Size score] $s_\text{size} = \min(r_A, r_B) / \max(r_A, r_B)$, where
        $r_A, r_B$ are the estimated radii (cube-root of voxel count). Penalizes
        large radius discrepancies.
\end{description}

The composite score is a weighted sum:
\begin{equation}
  s_\text{comp} = w_\text{prox} s_\text{prox} + w_\text{align} s_\text{align}
                + w_\text{cont} s_\text{cont} + w_\text{size} s_\text{size}
\end{equation}
with default weights $(0.35, 0.30, 0.25, 0.10)$ for the four components.

\paragraph{Distance-conditioned scoring.}
For long-range pairs (gap $> \texttt{long\_range\_threshold\_nm}$ = 1000~nm),
the proximity score decays below 0.05, which dominates the composite and
suppresses the score even when alignment and continuity are strong. To address
this, the pipeline switches to a proximity-free weight vector
$(0.00, 0.45, 0.40, 0.15)$ for long-range pairs, emphasizing alignment and
continuity. This was introduced in Experiment~14 and lifted ground-truth
coverage from 81.7\% to 83.5\% by recovering 27 long-range candidates that
had previously been suppressed by the proximity-dominated composite score
(see Chapter~\ref{ch:results} for details).

\section{Stage 4: Conservative Validation}

\paragraph{Three-outcome model.}
Every candidate is assigned one of three outcomes: ACCEPT (high confidence the
merge is correct), REJECT (high confidence it is not), or AMBIGUOUS (insufficient
evidence). Ambiguous candidates are preserved in the output and excluded from
precision/recall computation; they represent explicit uncertainty, not errors.

\paragraph{Validation rules.}
Seven rules are applied in sequence; any REJECT terminates evaluation of
that candidate:

\begin{enumerate}
  \item \textbf{MinGapRule}: Rejects candidates with gap $< \texttt{min\_gap\_nm}$.
        Intended to filter adjacency false positives (different axons sharing a
        voxel boundary) but also rejects genuine short-gap splits (Experiment~22
        shows this rule causes 37 false negatives in the current config).

  \item \textbf{MaxDistanceRule}: Hard-rejects if the endpoint-to-endpoint gap
        exceeds \texttt{max\_distance\_nm} (set to 20{,}000~nm for XPRESS to
        accommodate PCA-based endpoint pairs, whose tip-to-tip distance can reach
        12{,}000~nm).

  \item \textbf{CurvatureRule}: Rejects if the junction angle exceeds
        \texttt{max\_curvature\_deg} (150°). For long-range pairs
        (gap $> \texttt{skip\_distance\_nm}$ = 1000~nm), the rule is skipped
        entirely, because the endpoint-centroid direction estimate is unreliable
        at large distances and causes false rejections.

  \item \textbf{DirectionReversalRule}: Rejects if the two fragment tangent
        vectors point away from each other (dot product $< 0$).

  \item \textbf{SizeDiscrepancyRule}: Rejects if the cube-root radius ratio
        exceeds \texttt{max\_radius\_ratio} (15.0 for XPRESS).

  \item \textbf{BranchingLimitRule}: Rejects if the merge would connect to more
        than \texttt{max\_branches} (5{,}000) already-accepted partners. This is
        a safeguard against pathological graph nodes.

  \item \textbf{CompositeScoreRule}: Hard-rejects if the composite score is below
        \texttt{reject\_threshold} (0.15).
\end{enumerate}

All rules are configurable via YAML and may be omitted or reordered.

\paragraph{Accept/reject thresholds.}
After all rules pass, a candidate is accepted if its composite score exceeds
\texttt{accept\_threshold} (0.25). Candidates between \texttt{reject\_threshold}
and \texttt{accept\_threshold} are marked AMBIGUOUS.

\section{Stage 5: Assembly}

Accepted connections are added as edges to the fragment graph. Connected
components of this graph define assembled structures (reconstructed neurons
or axon segments). The assembler checks each structure for cycles, excessive
branching, and ambiguous connections, attaching topology warnings as metadata.
Per-structure confidence is computed as the minimum composite score across all
accepted connections (weakest-link model).

\section{Stage 6: Export}

The pipeline exports results in up to five formats: CSV tables of fragment
metadata and per-connection decisions (default, used for all XPRESS experiments),
GraphML (disabled by default due to OOM risk at scale, as 815{,}000+ edges produce
multi-gigabyte XML files), SWC neuron morphology, Neuroglancer annotation layers,
and a corrected precomputed segmentation volume with accepted merges applied.

\section{Ground Truth Evaluation}

\paragraph{The input--ground-truth contract.}
Before describing the two benchmarks, we state explicitly what the pipeline's
input is, what the ground truth is, and how the two are compared --- because
the comparison is not a direct segmentation-to-segmentation overlap.
\emph{Input}: a voxel-wise segmentation volume, the output of an automated
segmenter with known split errors.
\emph{Ground truth}: a set of labelled neurons that the pipeline's output is
scored against, provided in two different forms depending on the benchmark
(human-annotated segment labels for CREMI; manually traced skeleton graphs
for XPRESS).
\emph{Comparison}: the ground truth induces a set of fragment-pair labels
$\{(A, B) : \text{fragments } A \text{ and } B \text{ should be merged}\}$.
Each pipeline-accepted candidate $(A, B)$ is scored as a true positive if
the pair is in this ground-truth merge set, and as a false positive otherwise.
The pipeline's job is to propose exactly this set of pairs.

\paragraph{CREMI ground truth.}
For CREMI, the ground truth is available as human-annotated segment labels
on the raw volume. Two fragments $A$ and $B$ in the pipeline's output
constitute a ground-truth merge pair if they share the same human-annotated
label ID (i.e., they are pieces of the same annotated neuron split by the
pipeline's fragment extraction). Evaluation computes precision, recall, and
F1 over accepted candidates, excluding ambiguous candidates from all
denominators.

\paragraph{XPRESS skeleton ground truth.}
For XPRESS, ground truth is provided as manually traced skeleton graphs
(\texttt{XPRESS\_training\_skels.npz}) --- one graph per real axon, with each
node carrying a physical-nm coordinate. The ground-truth merge set is
constructed by an \emph{edge-crossing} rule: for every skeleton edge $(u, v)$,
we look up the segmentation label at $u$'s voxel and at $v$'s voxel in the
input segmentation; if the two endpoints land in different non-background
segment IDs $(\ell_u, \ell_v)$, then $(\ell_u, \ell_v)$ is added to the
ground-truth merge set. In words: when a real axon traced by the annotator
passes through two different baseline segments, those segments should be
merged, because they are pieces of the same neuron that the segmenter has
split. The ground-truth merge set is exactly the set of such pairs.

This construction (in \texttt{build\_merge\_oracle()}, named for historical
consistency with the XPRESS reference implementation) is conservative:
a pair enters the ground-truth set only when a skeleton edge physically
crosses a segment boundary, not merely when two segments contain skeleton
nodes from the same axon.

Two metrics are reported: \emph{Coverage} (fraction of ground-truth pairs for
which the pipeline generates at least one candidate), measuring
graph-construction reach; and \emph{Recall} (fraction of covered pairs that
are accepted), measuring validation accuracy. Precision is the fraction of
accepted candidates that correspond to ground-truth pairs.

\section{Visual Validation Report}
\label{sec:visual-validation}

To support qualitative inspection of pipeline behavior, we developed a diagnostic
report generator (\texttt{scripts/generate\_visual\_report.py}) that produces a
multi-page PDF or inline GitHub images.  For a stratified sample of candidates
across six categories --- \emph{ground-truth true positive} (accepted \emph{and}
ground-truth-positive), high-confidence accepted, low-confidence accepted,
borderline rejected, ground-truth false negative, and random rejected --- each
candidate occupies one row of three panels side-by-side:

\begin{enumerate}
  \item \textbf{Raw / Neutral:} Grayscale EM image (when available) or a
        neutral segmentation view in which each segment receives a distinct
        muted color but neither fragment is highlighted. Ground-truth skeleton
        nodes and edges near the Z-slice are overlaid as yellow dots/lines,
        showing where the manually traced neuron paths run through the tissue.
        This panel provides tissue context without the pipeline's coloring
        choices.

  \item \textbf{Pipeline Prediction:} Colorized segmentation with Fragment~A
        in red and Fragment~B in blue. A white arrow indicates the proposed
        connection direction; an ACC/REJ badge (green/red) states the decision.
        All five component scores (proximity, alignment, continuity, size,
        composite) and the gap distance are shown in the panel title.

  \item \textbf{Ground Truth:} Same colorized segmentation overlaid with the
        ground-truth verdict.  When a skeleton file is provided, skeleton
        nodes belonging to the axons passing through Fragments~A or~B are
        highlighted in bright yellow, and a dashed line marks the
        ground-truth edge crossing for positive pairs.  A large badge labels
        the outcome: \textbf{TP} (correct accept), \textbf{FP} (false
        accept), \textbf{FN} (missed merge), or \textbf{TN} (correct reject),
        together with an explicit ``Ground truth: SHOULD MERGE'' or
        ``Ground truth: NO MERGE'' annotation.
\end{enumerate}

This layout gives a human inspector all information needed to assess whether
the pipeline's decision is correct: the tissue context (Panel~1), the
pipeline's geometric reasoning (Panel~2), and the ground-truth verdict with
skeleton evidence (Panel~3).  Same-label candidates (two fragments from the
same segment, which trivially score 1.0) are excluded from sampling to avoid
contaminating the display with degenerate cases.

\paragraph{Ground-truth TP as the primary scientific check.}
The \emph{ground-truth TP} sampling category filters to the intersection of
pipeline-accepted and ground-truth-positive candidates and samples the top
rows by composite score. This is the only category that directly confirms
that the pipeline's acceptance decisions correspond to real biological
splits, rather than matching a metric count. For each ground-truth TP row,
the inspector verifies that (1) the yellow skeleton in Panel~1 runs through
both fragment regions, (2) the fragments look morphologically continuous in
Panel~2, and (3) the skeleton visibly occupies both label regions in Panel~3
with the dashed ground-truth-positive indicator.  The pipeline's best
configuration (Experiment~24) was inspected end-to-end using this category
and produced a dedicated PDF (\texttt{docs/visual\_report\_tp\_inspection.pdf})
that is committed to the repository for reproducibility.

\paragraph{CSV companion file.}
Alongside the PDF, the report writes a \texttt{gt\_pair\_audit.csv} file
containing one row per sampled candidate with the following columns:
\texttt{candidate\_id}, \texttt{fragment\_a}, \texttt{fragment\_b},
\texttt{label\_a}, \texttt{label\_b}, \texttt{gt\_should\_merge},
\texttt{gt\_source}, \texttt{status}, \texttt{composite}, \texttt{alignment},
\texttt{continuity}, and \texttt{gap}.  This enables programmatic downstream
analysis, for example cross-referencing a visual anomaly with the underlying
numerical scores or filtering by ground-truth agreement.
